\section{Riesgos del producto: un mundo de bajo daño también puede convertirse en uno de alto daño}

La visión de Elys es un mundo de bajo daño: menos fricción y conexiones más acertadas. Sin embargo, cualquier producto social puede causar daño, y la socialización mediada por IA amplifica ese riesgo. El doble digital puede representar mal al usuario, el context puede filtrarse, la IA puede fabricar una actividad social falsa y el algoritmo puede empujar al usuario hacia una exposición excesiva o hacia relaciones poco sanas.

\subsection{Lista de riesgos}

\noindent\begin{tabularx}{\textwidth}{p{0.22\textwidth}p{0.34\textwidth}X}
\toprule
Riesgo & Manifestación & Mecanismo necesario \\
\midrule
Representación errónea & El doble digital dice cosas que el usuario no respalda & Vista previa, posibilidad de retractarse y confirmación de acciones sensibles. \\
Filtración de privacidad & La plataforma u otras personas hacen mal uso del context de alta dimensionalidad & Almacenamiento local, cifrado y autorización por niveles. \\
Socialización falsa & Las interacciones con IA ahogan las interacciones entre personas reales & Señalar qué es IA y qué es una persona real; optimizar la tasa de conexión entre personas reales. \\
Manipulación y adicción & El doble digital genera dependencia emocional para aumentar la retención & Objetivos transparentes y límites de uso saludable. \\
Emparejamiento erróneo & El emparejamiento de alta dimensionalidad produce una falsa sensación de intimidad & Confirmación gradual, retroalimentación de personas reales y avisos sobre límites. \\
\bottomrule
\end{tabularx}

\begin{warningbox}{Un mundo de bajo daño no es un resultado automático}
La IA reduce la fricción, pero también puede reducir las defensas. El producto debe incorporar desde el diseño los permisos, la transparencia, la confirmación por parte de personas reales y los límites de uso saludable; de lo contrario, el mundo de bajo daño puede convertirse en un mundo de alto daño más difícil de detectar.
\end{warningbox}

\subsection{Resumen de la sección}

Los riesgos y el valor de la socialización mediada por IA provienen de lo mismo: te entiende mejor y también actúa con más iniciativa. Solo al comprender esto se pueden diseñar productos que tengan capacidad de conexión y, a la vez, sentido de los límites.
